% GENERATED FILE. Edit canonical lessons, then run npm run generate:books.
% canonical-source-hash: fnv1a64:5e5bc1f01a43b6b0
% canonical-lessons: RU-C463-priliv, RU-C463-otliv, RU-C463-bukhta, RU-C463-rakushka, RU-C463-nyryat

\chapter{At the seaside}
\label{ch:ru-at-the-seaside}
\hlchaptermodality{463}

\begin{chapteropening}
\textbf{By the end of this chapter:} \emph{I can describe the seaside: high tide, low tide, a bay, a seashell and diving.}

Complete the last lesson of chapter 463: I can describe the seaside: high tide, low tide, a bay, a seashell and diving.
\end{chapteropening}

\section[\texorpdfstring{\ru{прилив} (prilív)}{прилив (prilív)}]{\ru{прилив} (prilív) — a high tide, a flood tide; a surge}

\label{lesson:RU-C463-priliv}



\begin{quote}
Before the new one: say the Russian for to light, then the Russian for to put out.
\end{quote}

\subsection*{You'll want to know: \ru{прилив}}
\textbf{\ru{прилив}} (\emph{prilív}) — \textquotedblleft{}a high tide, a flood tide; a surge\textquotedblright{}.

\begin{quote}
\ru{Во} \ru{время} \ru{прилива} \ru{вода} \ru{подходит} \ru{к} \ru{дороге}.
\end{quote}

\begin{quote}
\emph{Vo vrémya prilíva vodá podkhódit k doróge.}
\end{quote}

\begin{quote}
At high tide the water comes up to the road.
\end{quote}

Masculine. \emph{\ru{Прилив} \ru{сил}} — a surge of energy.

\begin{cousinweb}
\emph{\ru{прилив}} is built from \emph{\ru{при}-} 'up to' and \emph{-\ru{лив}} 'a pouring', from \emph{\ru{лить}} 'to pour', which EtymDB takes from Proto-Slavic \emph{liti}.
\end{cousinweb}

\begin{grammarlens}[title={What you've built}]
One more word for telling, explaining and describing at B1.
\end{grammarlens}

\subsection*{Your turn}
\begin{itemize}
\raggedright
  \item \emph{Say it:} \emph{prilív}
  \item \emph{Say it:} \emph{prilív}, once more
  \item \emph{Say it:} a sentence of your own with \emph{prilív}
  \item \emph{Recall:} say the Russian for to light, then the Russian for to put out, then say \emph{prilív} again
\end{itemize}

\begin{tcolorbox}[breakable,colback=teal!4,colframe=teal!35!black,title={Before you move on}]
What does \ru{прилив} mean? (\textquotedblleft{}A high tide, a flood tide; a surge\textquotedblright{}.) Say it once more.
\end{tcolorbox}
\section[\texorpdfstring{\ru{отлив} (otlív)}{отлив (otlív)}]{\ru{отлив} (otlív) — a low tide, an ebb tide}

\label{lesson:RU-C463-otliv}



\begin{quote}
Before the new one: say the Russian for to put out, then the Russian for a high tide.
\end{quote}

\subsection*{You'll want to know: \ru{отлив}}
\textbf{\ru{отлив}} (\emph{otlív}) — \textquotedblleft{}a low tide, an ebb tide\textquotedblright{}.

\begin{quote}
\ru{Во} \ru{время} \ru{отлива} \ru{можно} \ru{дойти} \ru{до} \ru{острова} \ru{пешком}.
\end{quote}

\begin{quote}
\emph{Vo vrémya otlíva mózhno doytí do óstrova peshkóm.}
\end{quote}

\begin{quote}
At low tide you can walk to the island.
\end{quote}

Masculine.

\begin{cousinweb}
\emph{\ru{отлив}} is built from \emph{\ru{от}-} 'away' and \emph{-\ru{лив}} 'a pouring', from \emph{\ru{лить}} 'to pour', which EtymDB takes from Proto-Slavic \emph{liti}.
\end{cousinweb}

\begin{grammarlens}[title={What you've built}]
One more word for telling, explaining and describing at B1.
\end{grammarlens}

\subsection*{Your turn}
\begin{itemize}
\raggedright
  \item \emph{Say it:} \emph{otlív}
  \item \emph{Say it:} \emph{otlív}, once more
  \item \emph{Say it:} a sentence of your own with \emph{otlív}
  \item \emph{Recall:} say the Russian for to put out, then the Russian for a high tide, then say \emph{otlív} again
\end{itemize}

\begin{tcolorbox}[breakable,colback=teal!4,colframe=teal!35!black,title={Before you move on}]
What does \ru{отлив} mean? (\textquotedblleft{}A low tide, an ebb tide\textquotedblright{}.) Say it once more.
\end{tcolorbox}
\section[\texorpdfstring{\ru{бухта} (búkhta)}{бухта (búkhta)}]{\ru{бухта} (búkhta) — a bay, a cove}

\label{lesson:RU-C463-bukhta}



\begin{quote}
Before the new one: say the Russian for a high tide, then the Russian for a low tide.
\end{quote}

\subsection*{You'll want to know: \ru{бухта}}
\textbf{\ru{бухта}} (\emph{búkhta}) — \textquotedblleft{}a bay, a cove\textquotedblright{}.

\begin{quote}
\ru{Лодки} \ru{стоят} \ru{в} \ru{тихой} \ru{бухте}.
\end{quote}

\begin{quote}
\emph{Lódki stoyát v tíkhoy búkhte.}
\end{quote}

\begin{quote}
The boats are moored in a quiet bay.
\end{quote}

Feminine.

\begin{cousinweb}
\emph{\ru{бухта}} was borrowed from German \emph{Bucht} 'a bay, a bight' (EtymDB).
\end{cousinweb}

\begin{grammarlens}[title={What you've built}]
One more word for telling, explaining and describing at B1.
\end{grammarlens}

\subsection*{Your turn}
\begin{itemize}
\raggedright
  \item \emph{Say it:} \emph{búkhta}
  \item \emph{Say it:} \emph{búkhta}, once more
  \item \emph{Say it:} a sentence of your own with \emph{búkhta}
  \item \emph{Recall:} say the Russian for a high tide, then the Russian for a low tide, then say \emph{búkhta} again
\end{itemize}

\begin{tcolorbox}[breakable,colback=teal!4,colframe=teal!35!black,title={Before you move on}]
What does \ru{бухта} mean? (\textquotedblleft{}A bay, a cove\textquotedblright{}.) Say it once more.
\end{tcolorbox}
\section[\texorpdfstring{\ru{ракушка} (rakúshka)}{ракушка (rakúshka)}]{\ru{ракушка} (rakúshka) — a seashell}

\label{lesson:RU-C463-rakushka}



\begin{quote}
Before the new one: say the Russian for a low tide, then the Russian for a bay.
\end{quote}

\subsection*{You'll want to know: \ru{ракушка}}
\textbf{\ru{ракушка}} (\emph{rakúshka}) — \textquotedblleft{}a seashell\textquotedblright{}.

\begin{quote}
\ru{Дети} \ru{собирают} \ru{ракушки} \ru{на} \ru{пляже}.
\end{quote}

\begin{quote}
\emph{Déti sobiráyut rakúshki na plyázhe.}
\end{quote}

\begin{quote}
The children are collecting seashells on the beach.
\end{quote}

Feminine.

\begin{cousinweb}
\emph{\ru{ракушка}} is built on \emph{\ru{раковина}} 'a shell' (EtymDB), a word that Old East Slavic used for mother of pearl.
\end{cousinweb}

\begin{grammarlens}[title={What you've built}]
One more word for telling, explaining and describing at B1.
\end{grammarlens}

\subsection*{Your turn}
\begin{itemize}
\raggedright
  \item \emph{Say it:} \emph{rakúshka}
  \item \emph{Say it:} \emph{rakúshka}, once more
  \item \emph{Say it:} a sentence of your own with \emph{rakúshka}
  \item \emph{Recall:} say the Russian for a low tide, then the Russian for a bay, then say \emph{rakúshka} again
\end{itemize}

\begin{tcolorbox}[breakable,colback=teal!4,colframe=teal!35!black,title={Before you move on}]
What does \ru{ракушка} mean? (\textquotedblleft{}A seashell\textquotedblright{}.) Say it once more.
\end{tcolorbox}
\section[\texorpdfstring{\ru{нырять} (nyryát)}{нырять (nyryát)}]{\ru{нырять} (nyryát) — to dive}

\label{lesson:RU-C463-nyryat}



\begin{quote}
Before the new one: say the Russian for a bay, then the Russian for a seashell.
\end{quote}

\subsection*{You'll want to know: \ru{нырять}}
\textbf{\ru{нырять}} (\emph{nyryát}) — \textquotedblleft{}to dive\textquotedblright{}.

\begin{quote}
\ru{Мальчики} \ru{ныряют} \ru{с} \ru{камня} \ru{в} \ru{море}.
\end{quote}

\begin{quote}
\emph{Málchiki nyryáyut s kámnya v móre.}
\end{quote}

\begin{quote}
The boys dive from a rock into the sea.
\end{quote}

Imperfective; perfective \emph{\ru{нырнуть}}.

\begin{cousinweb}
\emph{\ru{нырять}} is usually linked with \emph{\ru{нора}} 'a burrow': to dive is to go down into a hole. EtymDB lists the word without a source, so the line is hedged.
\end{cousinweb}

\begin{grammarlens}[title={What you've built}]
One more word for telling, explaining and describing at B1.
\end{grammarlens}

\subsection*{Your turn}
\begin{itemize}
\raggedright
  \item \emph{Say it:} \emph{nyryát}
  \item \emph{Say it:} \emph{nyryát}, once more
  \item \emph{Say it:} a sentence of your own with \emph{nyryát}
  \item \emph{Recall:} say the Russian for a bay, then the Russian for a seashell, then say \emph{nyryát} again
\end{itemize}

\begin{tcolorbox}[breakable,colback=teal!4,colframe=teal!35!black,title={Before you move on}]
What does \ru{нырять} mean? (\textquotedblleft{}To dive\textquotedblright{}.) Say it once more.
\end{tcolorbox}
